\begingroup
\footnotesize
\renewcommand{\arraystretch}{1.16}
\setlength{\tabcolsep}{4pt}
\definecolor{mqQuestionBg}{gray}{0.93}
\definecolor{mqReady}{HTML}{1B7F3A}
\definecolor{mqDiagnostic}{HTML}{087F8C}
\definecolor{mqPolicy}{HTML}{6F42C1}
\definecolor{mqGate}{HTML}{D97706}
\definecolor{mqBlocked}{HTML}{B42318}
\newcommand{\mqStatus}[2]{\textcolor{#1}{\textbf{#2}}}
\newcommand{\mqPair}[4]{%
  \rowcolor{mqQuestionBg}%
  \multicolumn{2}{@{}L{0.96\textwidth}@{}}{\textbf{Question.} #3}\\*%
  \mqStatus{#1}{#2} & #4\\%
  \addlinespace[0.42em]%
}
\newcommand{\mqSeparator}[1]{%
  \addlinespace[0.5em]%
  \midrule
  \multicolumn{2}{@{}L{0.96\textwidth}@{}}{\textbf{#1}}\\*%
  \addlinespace[0.2em]%
}
\newcommand{\mqSubtable}[1]{%
  \par\Needspace{16\baselineskip}\medskip\noindent\textbf{#1}\par\smallskip%
}

\captionof{table}{Measurement-facing questions and current answers.  The rows
distinguish ready residuals, diagnostics, policy decisions, gates, and blocked
streams.  In each stream, the original current-answer rows are followed by a
separated block of additional single-topic questions needed for OGS grounding.
Where the project-specific answer is still unavailable, the row states the current
local/public evidence boundary and the missing provider or modelling decision.}
\label{tab:measurement-question-answers}

\noindent\textbf{Answer-status legend.}\par\smallskip
\begin{tabular}{@{}L{0.23\textwidth} L{0.70\textwidth}@{}}
\mqStatus{mqReady}{Ready} &
The local evidence and operator are sufficient for current active use.\\
\mqStatus{mqDiagnostic}{Diagnostic} &
The stream can be plotted or compared, but not yet weighted as an accepted
likelihood term.\\
\mqStatus{mqPolicy}{Policy} &
The answer is a current modelling choice, not a measurement fact.\\
\mqStatus{mqGate}{Gate} &
The stream is partly understood, but a collaborator decision, source file, or
uncertainty model is still required.\\
\mqStatus{mqBlocked}{Blocked} &
The local and public evidence do not provide the numeric export needed for a hard
residual.\\
\end{tabular}

\mqSubtable{Geometry And Stream Inventory}
\begin{longtable}{@{}L{0.18\textwidth} L{0.76\textwidth}@{}}
\toprule
\mqPair{mqReady}{Ready}
{Which local evidence layer defines the measurement positions?}
{The current geometry layer is the processed coordinate and mesh-lookup output:
\path{measurement_mesh_lookup.csv}, \path{borehole_mesh_lookup.csv}, and
\path{borehole_line_mesh_samples.csv}.  It transforms catalogue coordinates into
the two-dimensional OGS frame and records inside-cell, nearest-cell, and
outside-domain status.}
\mqPair{mqGate}{Gate}
{Can every measurement be treated as a point in the OGS mesh?}
{No.  NMR and RH are sparse supports, permeability is an interval-scale borehole
observation, Taupe/TDR is an EDZ-band line support, ERT is an inverted mesh field,
faults are 3D structural objects, and HM streams have source-frame geometry that
still needs a model-frame projection.}
\mqPair{mqReady}{Ready}
{Which streams have locally available measured-data plots?}
{Permeability, NMR, ERT, Taupe/TDR, RH, mini-piezometer pressure, crackmeter, and
levelling have useful measured-data or source-trend visualisations.  Generic
source-frame support plots for extensometer, laser, miniprisma, convergence, and
evapometer are catalogue evidence only; they are not useful paper figures until
values, times, units, and accepted OGS supports are attached.}
\mqSeparator{Additional granular questions for OGS grounding}
\mqPair{mqGate}{Gate}
{What raw source file defines this stream?}
{No single file defines all streams.  The authoritative raw sources are the
stream-specific files copied under \path{cda_knowledge_base/measurements_info/*/source_files}:
for example \path{Permeability_CDA_all_2025.xlsx}, weekly/seasonal NMR
\texttt{.dat} files, \path{ERT_meas_Niche_open.zip}, \path{Taupe_WC.xlsx},
\path{OT_RH5.xlsx}--\path{OT_RH8.xlsx}, and \path{VisualisationCDA.dat}.}
\mqPair{mqGate}{Gate}
{What processing step produced the table used here?}
{The tables used here come from the processed-observation/audit layer in
\path{SOTA_OGS_Mont_Terri_work/inversion_workflow/processed_observations} and the
joined \path{measurement_model_entry_matrix.md}.}
\mqPair{mqGate}{Gate}
{What exact spatial support does one row represent?}
{It is stream-specific: permeability rows are borehole intervals; NMR rows are
mapped NMR labels or profile positions; ERT rows are inverted mesh cells; Taupe rows
are sensor--EDZ-band line supports; RH rows are boundary-audit sensor points; HM rows
may be points, lines, pairs, or surfaces.}
\mqPair{mqGate}{Gate}
{What timestamp or survey epoch does one row represent?}
{It is stream-specific and must stay row-level: ERT/RH/Taupe/NMR carry dates or
timesteps, permeability rows are campaign/interval observations, and levelling rows
are survey-epoch differences.  Other HM time-series epochs are not locally exported.}
\mqPair{mqGate}{Gate}
{What unit and reference convention does the stored value use?}
{Known for active streams only: permeability is \(\mathrm{m^2}\) in log10 residual
space; NMR is vol.\% converted to fraction; ERT is resistivity/log-resistivity;
RH is percent converted to gauge liquid-pressure candidates.  Taupe and several HM
streams still need unit/reference confirmation.}
\mqPair{mqGate}{Gate}
{What uncertainty or covariance should be assigned to the row?}
{Only provisional scales are available: permeability uses 0.5 log10 units; NMR uses
reported confidence intervals with a 0.01 fraction floor plus bound-water caveats.
ERT, Taupe, RH boundary forcing, pressure, cracks, faults, and HM streams still need
accepted covariance or grouping rules.}
\mqPair{mqGate}{Gate}
{What OGS object does the stream ground?}
{The stream must be classified as an OGS input, output residual, material prior,
structural scenario, support mask, quality filter, or caveat before it can receive
weight.}
\bottomrule
\end{longtable}

\mqSubtable{Direct Permeability Pulse Tests}
\begin{longtable}{@{}L{0.18\textwidth} L{0.76\textwidth}@{}}
\toprule
\mqPair{mqReady}{Ready}
{What does the permeability stream measure?}
{It records interpreted nitrogen pulse-test permeability and transmissibility, plus
raw pressure-decay traces.  The active value is a noisy scalar interval
observation of intrinsic permeability in \(\mathrm{m^2}\), not hydraulic
conductivity, saturation, or liquid relative permeability.}
\mqPair{mqReady}{Ready}
{Where and when is it used in the current OGS objective?}
{The current active set is 75 positive rows on mapped BCD-A32 and BCD-A33 supports
inside the Niche 4 mesh.  They are used as 2024 quasi-static material observations
of \(\KK(\xx)\), with the older or outside-domain rows retained as context.}
\mqPair{mqPolicy}{Policy}
{How is the permeability residual formed?}
{The current residual is
\(\log_{10} k_{\mathrm{pred}}-\log_{10} k_{\mathrm{obs}}\), where
\(k_{\mathrm{pred}}\) is an interval-weighted directional response
\(\mathbf e^\top\KK\mathbf e\).  The first-pass standard deviation is
0.5 log10 units with duplicate-aware weights.}
\mqSeparator{Additional granular questions for OGS grounding}
\mqPair{mqGate}{Gate}
{Is each permeability value intrinsic liquid-equivalent permeability?}
{The active value is treated as an intrinsic-permeability candidate in
\(\mathrm{m^2}\), but the local workbooks and slides are gas pulse tests.  The
liquid-equivalent/slip-corrected interpretation is therefore a tracked caveat, not a
closed fact.}
\mqPair{mqGate}{Gate}
{Was a Klinkenberg or gas-slip correction applied?}
{Not documented row-by-row in the local workbooks.  Because nitrogen gas testing is
pressure dependent, the correction/provenance remains a gate for final
liquid-equivalent interpretation.}
\mqPair{mqGate}{Gate}
{What gas type and mean pressure were used for each test?}
{The local HERMES/method slides state nitrogen injection up to about 1 bar.  A
row-specific mean pressure for each interpreted permeability value is not extracted
from the current table.}
\mqPair{mqGate}{Gate}
{What temperature assumption was used in the gas-test interpretation?}
{Not found in the current processed table or source summaries.  Keep as a
gas-interpretation metadata gate unless the original fitting sheet or method note
states it.}
\mqPair{mqGate}{Gate}
{Which pressure-decay fits are failed or low-signal?}
{The processed audit separates 6687 pressure-decay rows and flags 4 interpreted rows
as missing or nonpositive.  It does not yet provide a row-level failed/low-signal fit
classification.}
\mqPair{mqGate}{Gate}
{Which rows may be affected by leakage?}
{No row-level leakage flag was found in the current local catalogue or audit layer.
Treat leakage as a requested provider/fitting-metadata flag, not as an inferred local
classification.}
\mqPair{mqGate}{Gate}
{What interval centre and interval length belong to each row?}
{Depth/interval centres are present in the interpreted tables; the source method uses
a nominal 10 cm static interval.  Test-specific interval length deviations, if any,
are not documented in the active target table.}
\mqPair{mqGate}{Gate}
{What borehole trace and direction belong to each row?}
{Available for current active rows: BCD-A32 and BCD-A33 have mapped endpoint geometry
and horizontal/vertical direction labels.  BCD-A24--A27 and BFM-D19 lack accepted
endpoint traces.}
\mqPair{mqGate}{Gate}
{What 3D volume is sampled by one pulse-test row?}
{Only approximately known: a 3D rock volume around the nominal 10 cm borehole interval
and its pressure-transient sensitivity.  The current OGS operator reduces this to a
2D interval-weighted directional average.}
\mqPair{mqGate}{Gate}
{Is the directional interval operator accepted for final scoring?}
{It is accepted for the current first-pass active objective, but final scoring still
needs a modelling-team policy decision on support conflicts, robust weighting, and
whether older endpoint-missing rows enter.}
\mqPair{mqGate}{Gate}
{Which bedding angle is used for permeability anisotropy?}
{The structural prior records a bedding angle near 144--145 degrees, but the table
does not yet prove whether this angle is fixed globally, released, or identical to
the mechanical/ERT anisotropy angle in the final workflow.}
\mqPair{mqGate}{Gate}
{Which repeated tests sample the same support?}
{The likelihood-policy audit groups repeated/conflicting active rows by support cell:
30 active repeated/range cells are audited, with 16 support groups showing range
\(\ge 2\) log10 units.}
\mqPair{mqGate}{Gate}
{How are several-log10 repeated-test disagreements classified?}
{They are currently tracked as support-conflict/outlier-policy evidence.  The active
row-Gaussian objective remains unchanged until a robust, support-cell, or
scalar-outlier policy is explicitly accepted.}
\mqPair{mqGate}{Gate}
{Are 2021 and 2024 rows both material observations?}
{They are both material/context observations in the catalogue, but only the 2024
mapped BCD-A32/A33 open-twin rows enter the current active objective.  2021 rows are
inactive unless their support and policy gates are reopened.}
\mqPair{mqGate}{Gate}
{How are older rows without endpoint geometry handled?}
{BCD-A24, BCD-A25, BCD-A26, BCD-A27, and BFM-D19 rows are retained visibly but
excluded as \texttt{missing\_segment\_geometry}; 98 positive rows would need labelled
endpoints or an accepted digitized trace before activation.}
\bottomrule
\end{longtable}

\mqSubtable{NMR Water Content}
\begin{longtable}{@{}L{0.18\textwidth} L{0.76\textwidth}@{}}
\toprule
\mqPair{mqDiagnostic}{Diagnostic}
{What physical quantity does NMR represent before modelling interpretation?}
{The local tables report volumetric water content in vol.\% and relaxation-time
information.  The signal is a total NMR-visible hydrogen/water signal, so the free
mobile water represented by \(\por\Sl\) is not separated from bound or interlayer
water by default.}
\mqPair{mqReady}{Ready}
{What are the available time ranges and supports?}
{The weekly 4E/4S tables have 170 rows from 2021-10-06 to 2025-09-09, and the
seasonal campaign archive has 294 profile rows.  The current OGS objective uses
the subset that maps to accepted Niche 4 support and output times.}
\mqPair{mqGate}{Gate}
{Why is raw absolute NMR water content not a final residual policy?}
{With fixed porosity \(n=0.105\), many rows exceed the maximum possible
\(\theta=\por\Sl\).  This shows that a bound/interlayer-water bias, campaign bias,
or trend/anomaly residual is needed before NMR can be interpreted as mobile OGS
water.}
\mqPair{mqPolicy}{Policy}
{What is the preferred next NMR residual?}
{The preferred candidate is a within-label trend/anomaly residual.  It compares
temporal changes after subtracting label means, so constant bound-water or campaign
offsets cancel to first order.  Absolute or bias-corrected water-content residuals
remain gated until the mobile fraction and uncertainty policy are accepted.}
\mqSeparator{Additional granular questions for OGS grounding}
\mqPair{mqGate}{Gate}
{Which NMR device produced each row?}
{Weekly rows are from the 4S/4E single-sided NMR monitoring files; seasonal rows are
from the Niche 3/Niche 4 seasonal campaign archive.  Row-level device metadata beyond
those labels is not expanded in the current table.}
\mqPair{mqGate}{Gate}
{Which campaign produced each row?}
{Weekly rows come from \path{Weekly_2021-2022_at_4S.dat} and
\path{Weekly_2022-2025_at_4E.dat}; seasonal rows are dated files named
\path{Niche*_Y_YYYY_M_MM_D_DD.dat}.}
\mqPair{mqGate}{Gate}
{What borehole label and depth interval does each row sample?}
{The processed NMR targets keep label/profile position fields and are mapped through
the NMR/Taupe coordinate workbook.  The paper table does not repeat every row-level
label/depth; those details live in \path{nmr_weekly.csv} and
\path{nmr_seasonal_profiles.csv}.}
\mqPair{mqGate}{Gate}
{What sensitive footprint does one NMR row represent?}
{Only approximately local support is represented in the current operator.  The exact
single-sided NMR sensitive volume or profile footprint is not documented row-by-row,
so current comparisons sample mapped labels rather than a full sensitivity kernel.}
\mqPair{mqGate}{Gate}
{Which fitting model produced the reported water content?}
{Not fully reconstructed from the local processed tables.  The local NMR
presentations discuss relaxation-time and bound/interlayer-water issues, but the
row-level inversion/fitting model is still a source-method gate.}
\mqPair{mqGate}{Gate}
{Which labels are accepted as inside the Niche 4 OGS support?}
{The processed layer maps 287 of 464 NMR targets to usable current-mesh support and
uses 192 rows in the active state objective.  Niche 3 seasonal rows remain outside
the current Niche 4 fit.}
\mqPair{mqGate}{Gate}
{Are NMR dates acquisition dates or processed-table dates?}
{Seasonal file names encode campaign acquisition dates; weekly files carry monitoring
dates.  No separate processing-date convention is used for OGS timing in the current
processed targets.}
\mqPair{mqGate}{Gate}
{How should NMR confidence intervals become standard deviations?}
{The current likelihood converts reported 95 percent confidence intervals to a row
scale and applies a 0.01 fraction uncertainty floor.  This handles measurement spread,
not the unresolved bound-water bias.}
\mqPair{mqGate}{Gate}
{Does the stored NMR value represent total NMR-visible water?}
{Yes for current modelling purposes: it is treated as total NMR-visible water content,
which can include mobile plus bound/interlayer contributions.}
\mqPair{mqGate}{Gate}
{Does the stored NMR value represent mobile pore water?}
{Not by default.  Mobile OGS water is approximated by \(\por\Sl\); the stored NMR
value needs a bias/free-water/trend policy before being interpreted as only mobile
pore water.}
\mqPair{mqGate}{Gate}
{Can bound water be treated as constant by label?}
{This is the preferred first simplification for trend/anomaly residuals: subtract
within-label means so constant label offsets cancel to first order.  It remains a
policy choice, not a measured decomposition.}
\mqPair{mqGate}{Gate}
{Can bound water be treated as constant by campaign or depth?}
{Not proven.  The bound-water audit shows high-offset labels and nonuniform offsets,
so a single campaign/depth correction is riskier than label-specific trend/anomaly or
bias terms.}
\mqPair{mqGate}{Gate}
{Is a T2 threshold available for the mobile fraction?}
{The NMR presentation discusses a possible cut-off near 0.5--0.6 ms, but no accepted
row-level T2 threshold has been promoted to the OGS residual policy.}
\mqPair{mqGate}{Gate}
{Is there a drying, CCM, or other calibration for bound water?}
{Public/local CD-A sources discuss NMR, CCM, and water-content interpretation, but no
accepted calibration has been converted into a bound-water subtraction for these
targets.}
\mqPair{mqGate}{Gate}
{How should rows with \(\theta_{\mathrm{obs}}>n\) be used?}
{They should not be used as raw absolute mobile-water residuals until the
bound-water or bias policy is accepted.}
\bottomrule
\end{longtable}

\mqSubtable{ERT Resistivity}
\begin{longtable}{@{}L{0.18\textwidth} L{0.76\textwidth}@{}}
\toprule
\mqPair{mqDiagnostic}{Diagnostic}
{What does ERT measure?}
{ERT measures an electrical response that is inverted to a resistivity field.  It
does not directly measure saturation, water content, pressure, or permeability.}
\mqPair{mqDiagnostic}{Diagnostic}
{How should OGS be tied to ERT?}
{The forward direction is
\(\Sl\rightarrow\theta=\por\Sl\rightarrow\rho_{\Omega\mathrm{m}}\rightarrow
\mathcal O_{\mathrm{ERT}}\).  The current local default uses
\(\rho_{\Omega\mathrm{m}}=1.108\,\theta^{-1.58}\) and compares
\(\log_{10}\rho\), not raw water content.}
\mqSeparator{Additional granular questions for OGS grounding}
\mqPair{mqGate}{Gate}
{Is the ERT residual target resistivity?}
{The physical observable is inverted bulk resistivity.  The current residual should
use resistivity after transforming OGS \(\theta=\por\Sl\) through the selected
water--resistivity law.}
\mqPair{mqGate}{Gate}
{Is the ERT residual target log-resistivity?}
{The current diagnostic comparison uses \(\log_{10}\rho\), but final residual use
still needs accepted support and covariance.}
\mqPair{mqGate}{Gate}
{Is the ERT residual target resistivity change?}
{Not in the current diagnostic.  A change/anomaly target is possible but would need a
reference date and covariance policy that are not accepted yet.}
\mqPair{mqGate}{Gate}
{Which inversion mesh is authoritative?}
{The current projection uses the reference VTK mesh from
\path{Niche open/2019/11-2019/dcinv.result_01.vtk}, with 5966 cells in the
projection lookup.}
\mqPair{mqGate}{Gate}
{Which VTK field is authoritative?}
{The processed audit uses the inverted \texttt{Resistivity(log10)} field from the
ERT VTK files for diagnostic comparisons.}
\mqPair{mqGate}{Gate}
{Which reference date is used for ERT changes or anomalies?}
{None is accepted for the current residual; the current diagnostic compares
log10-resistivity fields, not relative changes against a selected baseline.}
\mqPair{mqGate}{Gate}
{Are bad electrodes or bad timesteps flagged?}
{Timesteps are inventoried and 16 rows lack matching VTK files.  A detailed
bad-electrode/bad-timestep quality table is not locally accepted for residual
weighting.}
\mqPair{mqGate}{Gate}
{Are inversion artefacts flagged?}
{No residual-ready artefact mask or inversion-quality covariance is available.  This
is one reason ERT stays diagnostic.}
\mqPair{mqGate}{Gate}
{Which water-resistivity law should be used?}
{The current local diagnostic uses \(\rho_{\Omega\mathrm{m}}=1.108\,\theta^{-1.58}\);
final acceptance is still gated.}
\mqPair{mqGate}{Gate}
{What uncertainty is assigned to the water-resistivity exponent?}
{No accepted exponent uncertainty is recorded.  Alternative fitted relations are
retained for sensitivity, with the Niche-4 paired diagnostic RMSE about 0.202 log10
resistivity units.}
\mqPair{mqGate}{Gate}
{What pore-fluid resistivity should be used?}
{No separate pore-fluid resistivity is used in the current empirical transform; it is
lumped into \(\rho=1.108\,\theta^{-1.58}\).  A physics-based Archie law would require
this value.}
\mqPair{mqGate}{Gate}
{What porosity should be used in the ERT transform?}
{The current diagnostic uses the OGS porosity field through
\(\theta=\por\Sl\); in the frozen current model this is the fixed porosity support
used elsewhere in the state sampling.}
\mqPair{mqGate}{Gate}
{Does the transform use total water content \(\theta\)?}
{The current diagnostic transform uses \(\theta=\por\Sl\).}
\mqPair{mqGate}{Gate}
{Does the transform need a free-water correction?}
{Unresolved.  The current diagnostic uses total model \(\theta=\por\Sl\), but clay
bound/interlayer water and surface conduction may require an observation-model
correction before final use.}
\mqPair{mqGate}{Gate}
{Does clay surface conductivity need an additive term?}
{Possibly.  Clay/shaly-sand literature and the water-content paper identify
surface/interfacial conductivity as a caveat, but no CD-A additive term is accepted
in the current operator.}
\mqPair{mqGate}{Gate}
{Does bound water conduct like mobile pore water?}
{Not established.  The current empirical transform does not separate bound and mobile
water, so this stays an ERT/NMR calibration gate.}
\mqPair{mqGate}{Gate}
{Can the calibration separate porosity from water content?}
{No.  The current relation maps \(\theta\) to resistivity and does not independently
identify porosity and saturation.}
\mqPair{mqGate}{Gate}
{Can the calibration separate mineralogy from water content?}
{No.  Lithology/mineralogy effects are retained as interpretation caveats or future
masks, not separated by the present regression.}
\mqPair{mqGate}{Gate}
{Can the calibration separate pore-fluid chemistry from water content?}
{No.  Pore-fluid chemistry is not an independent variable in the current diagnostic
relation.}
\mqPair{mqGate}{Gate}
{Is the ERT inversion electrically isotropic?}
{Not confirmed from the local inversion metadata.  The current diagnostic treats the
inverted resistivity field as a scalar field.}
\mqPair{mqGate}{Gate}
{What resistivity tensor is used if anisotropy matters?}
{None in the current OGS--ERT diagnostic.  A tensor resistivity observation model
would require accepted electrical anisotropy data and orientation.}
\mqPair{mqGate}{Gate}
{Does the ERT bedding angle match the hydraulic bedding angle?}
{Unknown.  A local \path{bedding_angle_ERT.png} exists and the structural prior is
near 144--145 degrees, but equality with the hydraulic tensor angle is not accepted.}
\mqPair{mqGate}{Gate}
{Does the ERT bedding angle match the mechanical bedding angle?}
{Unknown for final scoring.  The report keeps bedding as a structural prior until
mechanical, hydraulic, and electrical anisotropy conventions are reconciled.}
\mqPair{mqGate}{Gate}
{Do fractures require a separate ERT mask or term?}
{Probably for final interpretation: public/local CD-A sources link high resistivity
regions with open fractures or preferential paths.  No accepted fracture mask is
active in the current residual.}
\mqPair{mqGate}{Gate}
{Do carbonate layers require a separate ERT mask or term?}
{Not found as an accepted model mask in the local operator.  Lithology/facies effects
remain structural caveats unless a mask is supplied.}
\mqPair{mqGate}{Gate}
{What ERT-to-OGS coordinate transform is accepted?}
{None is externally accepted yet.  The current reproducible provisional transform is
\(x_{\mathrm{model}}=x_{\mathrm{raw}}\), \(y_{\mathrm{model}}=y_{\mathrm{raw}}+500\).}
\mqPair{mqGate}{Gate}
{What near-niche ERT mask is accepted?}
{No final mask is accepted.  The current approximate support keeps 2035 cells that
are inside an OGS cell and inside the provisional near-niche support.}
\mqPair{mqGate}{Gate}
{Is the 35 cm rock-depth convention part of the ERT support?}
{It is an explicit open support question.  The current support-sensitivity audit
tests radial subsets, but the exact 35 cm convention is not closed.}
\mqPair{mqGate}{Gate}
{How are dense ERT cells aggregated before scoring?}
{They are not yet scored.  Future scoring must aggregate or weight by support/time so
correlated inversion cells are not treated as independent point data.}
\mqPair{mqGate}{Gate}
{What spatial covariance is assigned to ERT cells?}
{No defensible spatial covariance is recorded locally or in the current public
sources for these CD-A VTK cells.}
\mqPair{mqGate}{Gate}
{What temporal covariance is assigned to ERT timesteps?}
{No temporal covariance is accepted.  Time grouping and sampling tolerance remain
part of the ERT likelihood gate.}
\bottomrule
\end{longtable}

\mqSubtable{Taupe/TDR}
\begin{longtable}{@{}L{0.18\textwidth} L{0.76\textwidth}@{}}
\toprule
\mqPair{mqDiagnostic}{Diagnostic}
{What does the Taupe/TDR workbook currently provide?}
{It provides 5088 dated sensor-band rows from A3, A4, A7, and A8 over
2019-12-01 to 2025-10-10.  The workbook values are kept in source units until the
provider confirms whether they are calibrated volumetric water content, apparent
relative dielectric permittivity, ARDP, or another Taupe proxy.}
\mqPair{mqDiagnostic}{Diagnostic}
{Which Taupe supports can be mapped to the current OGS slice?}
{A3 and A4 are mapped line supports in the current mesh.  A7 and A8 are retained in
the data tables but are outside the current local OGS support.}
\mqPair{mqPolicy}{Policy}
{How can Taupe/TDR enter OGS before absolute calibration is known?}
{Use only standardized same-series trend or anomaly diagnostics by sensor and EDZ
band.  A hard absolute water-content or permittivity residual remains blocked.}
\mqSeparator{Additional granular questions for OGS grounding}
\mqPair{mqGate}{Gate}
{Does the workbook value store apparent relative dielectric permittivity?}
{Possible but not confirmed for \path{Taupe_WC.xlsx}.  Public/local TDR context says
TAUPE/TDR can measure apparent relative dielectric permittivity, but the workbook
unit has not been provider-confirmed.}
\mqPair{mqGate}{Gate}
{Does the workbook value store ARDP?}
{Possible.  The current semantic audit treats the values as an ARDP/dielectric
water-content proxy until the unit is confirmed.}
\mqPair{mqGate}{Gate}
{Does the workbook value store calibrated volumetric water content?}
{Not confirmed.  Interpreting the values as water-content percent is only a
plausibility check in the audit, not an accepted calibration.}
\mqPair{mqGate}{Gate}
{Does the workbook value store travel time or amplitude?}
{No evidence in the processed workbook summary.  The sheets contain dated band values
near 9.87--11.00, not raw waveform traces, travel times, or amplitudes.}
\mqPair{mqGate}{Gate}
{What sensor-specific calibration converts Taupe/TDR to water content?}
{Not present in the local catalogue.  This is the main blocker for absolute
Taupe/TDR water-content residuals.}
\mqPair{mqGate}{Gate}
{What Taupe/TDR calibration range is valid for Opalinus Clay?}
{Not accepted for the CD-A workbook values.  The audit only tests physical
plausibility of candidate conversions; it does not establish an Opalinus Clay
calibration range.}
\mqPair{mqGate}{Gate}
{Which electrode tube does each row represent?}
{Rows are grouped by sheet/sensor labels A3, A4, A7, and A8.  The exact electrode
tube geometry is mapped through the Taupe/NMR coordinate workbook and
\path{VisualisationCDA.dat}, but a row-level tube metadata table is not printed here.}
\mqPair{mqGate}{Gate}
{Which EDZ band or depth interval does each row represent?}
{Each sheet contains EDZ distance bands from the niche wall: 0--10, 10--20, 20--30,
30--40, 40--50, and aggregate 0--50 cm bands.}
\mqPair{mqGate}{Gate}
{Are A7 and A8 outside the current OGS mesh?}
{They are retained as outside-support context in the current local tables.}
\mqPair{mqGate}{Gate}
{Are air gaps or tube-contact problems flagged?}
{No row-level air-gap/contact-quality flags were found in the current Taupe/TDR
workbook or semantic audit.}
\mqPair{mqGate}{Gate}
{Is convergence or settling time flagged for Taupe/TDR?}
{No explicit convergence/settling flag is present in the processed Taupe table.
Temporal startup effects would need a provider or Geoscope quality flag.}
\mqPair{mqGate}{Gate}
{Are failed Taupe/TDR sensor segments flagged?}
{No accepted failed-segment flag is present.  Current gating is by support: A3/A4
mapped in-domain, A7/A8 outside the present mesh.}
\mqPair{mqGate}{Gate}
{Does bound water affect Taupe/TDR permittivity like mobile water?}
{Unknown.  Dielectric response can differ from mobile OGS water, especially in
claystone, so this remains part of the Taupe calibration/observation-model gate.}
\mqPair{mqGate}{Gate}
{Is a temperature correction needed for Taupe/TDR?}
{Not determined from the local workbook.  No temperature correction is applied in the
current trend diagnostic.}
\mqPair{mqGate}{Gate}
{Is a salinity or conductivity-loss correction needed for Taupe/TDR?}
{Not determined.  TDR can be affected by electrical conductivity, but the current
local Taupe operator has no salinity/conductivity-loss correction.}
\mqPair{mqGate}{Gate}
{Should Taupe/TDR be compared as absolute water content?}
{No for the current workflow.  Absolute comparison remains blocked until unit,
calibration, uncertainty, and bound-water/dielectric interpretation are accepted.}
\mqPair{mqGate}{Gate}
{Should Taupe/TDR be compared as a standardized anomaly?}
{This is the current safe diagnostic policy before absolute calibration is known.}
\bottomrule
\end{longtable}

\mqSubtable{RH, Suction, And Boundary Pressure}
\begin{longtable}{@{}L{0.18\textwidth} L{0.76\textwidth}@{}}
\toprule
\mqPair{mqDiagnostic}{Diagnostic}
{What does RH measure and how is it converted?}
{The RH workbooks measure vapour relative humidity in percent.  The local audit
converts it to a gauge liquid-pressure candidate with
\(\pl=\rho_lRT/M_w\,\ln(\mathrm{RH})\), assuming \(T=298.15\,\mathrm{K}\) and
\(\rho_l=1095\,\mathrm{kg\,m^{-3}}\).}
\mqPair{mqDiagnostic}{Diagnostic}
{Does RH tie to OGS input or output?}
{Primarily input.  RH/suction can audit or reconstruct the open-niche pressure
boundary \(p_E(t)\), and later can inform retention or boundary uncertainty.  It
is not an interior material residual by itself.}
\mqSeparator{Additional granular questions for OGS grounding}
\mqPair{mqGate}{Gate}
{Which RH sensor defines the open-niche boundary candidate?}
{No single sensor is accepted as the boundary source.  The audit prefers RH5/RH6 for
clean open-twin evidence and constructs a candidate RH5/RH6 median curve, but the
active \(p_E(t)\) curve provenance is not confirmed.}
\mqPair{mqGate}{Gate}
{Which temperature sensor should accompany the RH sensor?}
{Not identified in the active boundary curve provenance.  The local Kelvin audit uses
a fixed \(T=298.15\,\mathrm{K}\) rather than a row-level temperature series.}
\mqPair{mqGate}{Gate}
{Where is the RH sensor relative to \(\Gamma_E\)?}
{The source has \path{Location_suc.png} and mapped RH/suction support points near the
open niche.  The exact relation between each sensor and the model boundary
\(\Gamma_E\) is still a boundary-support gate.}
\mqPair{mqGate}{Gate}
{Are door openings filtered from the RH data?}
{Not in the local RH audit.  Door/opening times are mentioned as Geoscope auxiliary
data, but the numeric event export is absent.}
\mqPair{mqGate}{Gate}
{Are logger problems filtered from the RH data?}
{Partly by rule, not by provider log: RH7/RH8 low values below 50 percent are treated
as suspicious/outliers.  No complete logger-status table is present.}
\mqPair{mqGate}{Gate}
{Are high-RH saturation intervals filtered?}
{They are flagged with caution, not removed by a final policy.  Open-twin
thermo-hygrometer values above 95 percent RH are considered less reliable.}
\mqPair{mqGate}{Gate}
{Is RH stored as percent or fraction before conversion?}
{The workbooks store RH in percent.  The Kelvin conversion divides by 100 and uses RH
as a fraction inside the logarithm.}
\mqPair{mqGate}{Gate}
{What temperature is used in the Kelvin conversion?}
{The current audit uses \(T=298.15\,\mathrm{K}\); final boundary provenance remains
gated.}
\mqPair{mqGate}{Gate}
{What liquid density is used in the Kelvin conversion?}
{The current audit uses \(\rho_l=1095\,\mathrm{kg\,m^{-3}}\); final boundary
provenance remains gated.}
\mqPair{mqGate}{Gate}
{Does the RH conversion produce capillary pressure?}
{The audit computes a gauge liquid-pressure candidate.  Capillary pressure follows by
the sign convention \(p_c=-\pl\) if gas pressure is taken as the reference.}
\mqPair{mqGate}{Gate}
{Does the RH conversion produce gauge liquid pressure?}
{The current audit interprets the result as a gauge liquid-pressure candidate.}
\mqPair{mqGate}{Gate}
{Does the RH sign convention match \(p_c=-\pl\)?}
{Yes in the current interpretation: RH below 100 percent gives negative
gauge-liquid pressure, equivalent to positive suction/capillary pressure.  The active
curve provenance still needs confirmation.}
\mqPair{mqGate}{Gate}
{Was the current \(p_E(t)\) curve derived from RH data?}
{Not proven.  The active curve implies much drier RH than clean RH5/RH6 over most of
the overlap, so it must be treated as existing model forcing with unverified
provenance.}
\mqPair{mqGate}{Gate}
{Was the current \(p_E(t)\) curve manually constructed?}
{Unknown.  No local generation script, smoothing record, or provider note explains
whether it was manual, RH-derived, or based on another source.}
\mqPair{mqGate}{Gate}
{What smoothing was applied to \(p_E(t)\)?}
{Not documented locally.  This is included in the RH active-curve provenance request.}
\mqPair{mqGate}{Gate}
{What interpolation policy is used for \(p_E(t)\)?}
{Not documented as a source-data policy.  OGS can interpolate curves internally, but
the intended preprocessing and date-to-model-time convention remain open.}
\mqPair{mqGate}{Gate}
{What extrapolation policy is used outside the \(p_E(t)\) date range?}
{Not documented.  The active curve spans model dates 2019-09-18 to 2023-12-26, while
local RH data extend to 2025-09-04.}
\mqPair{mqGate}{Gate}
{What uncertainty is assigned to \(p_E(t)\)?}
{No accepted boundary-forcing uncertainty is assigned.  The RH candidate envelope has
p50 pressure spread about 2.10 MPa, but it is not an accepted likelihood model.}
\bottomrule
\end{longtable}

\mqSubtable{Direct Mini-Piezometer Pressure}
\begin{longtable}{@{}L{0.18\textwidth} L{0.76\textwidth}@{}}
\toprule
\mqPair{mqDiagnostic}{Diagnostic}
{What does the direct pressure stream measure?}
{It is the mini-piezometer stream.  The public CD-A inventory defines
mini-piezometer output as pore pressure in Pa.  The 2026 source slide plots
A28--A31 in kPa from about March 2024 to January 2026.}
\mqSeparator{Additional granular questions for OGS grounding}
\mqPair{mqGate}{Gate}
{What are the accepted OGS coordinates for BCD-A28--A31?}
{Not accepted yet.  The slide has a location inset and the layout catalogue has
support geometry, but no machine-readable BCD-A28--A31 coordinate/support table has
been promoted to OGS pressure targets.}
\mqPair{mqGate}{Gate}
{What screened interval belongs to each pressure sensor?}
{Not found locally.  The requested export must include screened interval, sensor
depth/elevation, or trace/support id for each mini-piezometer.}
\mqPair{mqGate}{Gate}
{What depth or elevation belongs to each pressure sensor?}
{Not found in residual-ready form.  Depth/elevation is required before comparing to
OGS liquid pressure or hydraulic head.}
\mqPair{mqGate}{Gate}
{What exact timestamps are available for the pressure stream?}
{Only source-plot timing is visible: roughly March 2024 to January 2026.  Exact
timestamps are absent from the local catalogue.}
\mqPair{mqGate}{Gate}
{What sampling interval is used by the pressure stream?}
{Not found.  The Geoscope export request asks for sampling interval and timestamp
convention.}
\mqPair{mqGate}{Gate}
{What maintenance flags exist for the pressure stream?}
{The May 2026 minutes state BCD-A28--A31 were working well, but no row-level
maintenance or quality flags are locally exported.}
\mqPair{mqGate}{Gate}
{Are the exported pressure values pore pressure?}
{The public inventory describes mini-piezometer output as pore pressure, but the
machine-readable export still needs confirmation.}
\mqPair{mqGate}{Gate}
{Are the exported pressure values gauge pressure?}
{Unknown.  The slide is in kPa and public inventory says pore pressure in Pa, but the
absolute/gauge/head convention is not confirmed.}
\mqPair{mqGate}{Gate}
{Are the exported pressure values hydraulic head?}
{Unknown.  Hydraulic-head use would require elevation/reference metadata that is not
present locally.}
\mqPair{mqGate}{Gate}
{What pressure unit is used in the raw export?}
{No raw export is present.  The local source slide plots kPa; the public CD-A
inventory lists mini-piezometer pore pressure in Pa.}
\mqPair{mqGate}{Gate}
{What atmospheric correction is applied?}
{Not documented.  This is part of the pressure reference-convention gate.}
\mqPair{mqGate}{Gate}
{What elevation correction is applied?}
{Not documented.  Required before converting between pressure and head or comparing
sensors at different elevations.}
\mqPair{mqGate}{Gate}
{Which pressure sensors are in the open twin?}
{The local trend slide distinguishes open/closed twin behaviour for A28--A31, but the
accepted sensor-to-twin mapping is not encoded in a processed target table.}
\mqPair{mqGate}{Gate}
{Which pressure sensors are in the closed twin?}
{Not encoded in a residual-ready table.  Must be taken from the future Geoscope export
or confirmed from the slide/support geometry.}
\mqPair{mqGate}{Gate}
{Which pressure sensors are far-field supports?}
{Not classified locally.  The current question remains whether any mini-piezometer
samples far-field pressure, niche boundary evolution, or a local 3D path.}
\mqPair{mqGate}{Gate}
{Which pressure sensors are outside the 2D slice?}
{Not determined because accepted coordinates/supports are missing.}
\mqPair{mqGate}{Gate}
{Does the pressure stream validate the far-field Dirichlet pressure?}
{Potentially, but not yet.  It needs sensor support, reference convention, and
uncertainty before it can test the far-square hydraulic Dirichlet boundary.}
\mqPair{mqGate}{Gate}
{Does the pressure stream validate the niche boundary evolution?}
{Potentially, especially if a sensor is near the open niche, but support and pressure
reference are not yet known.}
\mqPair{mqGate}{Gate}
{Does the pressure stream validate the initial pressure field?}
{Potentially if observations exist near model start or post-excavation start.  The
available local plot starts around 2024, so it cannot by itself reconstruct the
initial field.}
\mqPair{mqGate}{Gate}
{Can the 2D model represent the sampled hydraulic path?}
{Unknown.  This depends on whether each mini-piezometer samples a path close to the
2D slice or a 3D feature that the current model cannot represent.}
\bottomrule
\end{longtable}

\mqSubtable{Crackmeter And Visible Surface Cracks}
\begin{longtable}{@{}L{0.18\textwidth} L{0.76\textwidth}@{}}
\toprule
\mqPair{mqDiagnostic}{Diagnostic}
{What does the crackmeter plot measure?}
{The source slide reports post-processed deformation in mm from roughly August
2021 to January 2026 for N3/N4 crackmeter or extensometer-labelled crack supports.
It is trend evidence for local closure/opening and seasonal response, not yet a
hard displacement residual.}
\mqSeparator{Additional granular questions for OGS grounding}
\mqPair{mqGate}{Gate}
{Is the crackmeter value crack aperture?}
{Possible but not confirmed.  The slide calls the plotted quantity post-processed
deformation in mm; the export must state whether it is aperture, displacement, or an
instrument deformation.}
\mqPair{mqGate}{Gate}
{Is the crackmeter value relative displacement?}
{Possible.  Public seasonal-HM sources describe crackmeters as measuring deformation
perpendicular to a crack, but the CD-A source export must define the component.}
\mqPair{mqGate}{Gate}
{Is the crackmeter value strain?}
{No evidence that the plotted crackmeter value is strain.  The units shown are mm,
so displacement/aperture is more likely, pending confirmation.}
\mqPair{mqGate}{Gate}
{Does positive value mean opening or closure?}
{Not confirmed.  The sign convention and whether positive means opening or closure
must be supplied with the crackmeter export.}
\mqPair{mqGate}{Gate}
{What baseline date is used for the crackmeter trend?}
{Not documented in the table.  The plotted series spans roughly August 2021 to
January 2026, but the zero/reference date is absent.}
\mqPair{mqGate}{Gate}
{What instrument axis is measured by each crackmeter?}
{Not documented in residual-ready form.  Axis/orientation is required before mapping
to OGS displacement or strain.}
\mqPair{mqGate}{Gate}
{Which visible crack does each time series represent?}
{The slide has an instrument inset and the Tecplot layout contains
fracture/crack-geometry zones, but no accepted series-to-crack mapping table is
present.}
\mqPair{mqGate}{Gate}
{Should crack data be compared to displacement jump?}
{Only if a crack/fault discontinuity observation operator is introduced.  The current
2D continuum model has no direct displacement-jump residual for crackmeters.}
\mqPair{mqGate}{Gate}
{Should crack data be compared to near-niche strain?}
{Possible as a qualitative or projected diagnostic after axis, support, and sign are
known.  It is not currently an active residual.}
\mqPair{mqGate}{Gate}
{Can visible cracks define a local EDZ prior?}
{Yes as a plausible structural prior or mask, if the crack geometry is projected to
the 2D slice and not over-weighted as a time-series residual.}
\mqPair{mqGate}{Gate}
{Can visible cracks define a local permeability prior?}
{Yes, but only as prior/scenario evidence unless independent hydraulic evidence
supports a permeability anomaly.}
\mqPair{mqGate}{Gate}
{What uncertainty should be assigned to crackmeter values?}
{No instrument uncertainty or covariance is available locally.  Public sources
confirm the measurement type, not the CD-A residual weight.}
\bottomrule
\end{longtable}

\mqSubtable{Extensometer, Convergence, Miniprisma, Laser, And Levelling}
\begin{longtable}{@{}L{0.18\textwidth} L{0.76\textwidth}@{}}
\toprule
\mqPair{mqDiagnostic}{Diagnostic}
{What does the precision-levelling summary measure?}
{It gives vertical height differences in mm between the initial 2022-08-29/30
survey and the 2026-03 campaign at 12 survey points.  It can be compared to
vertical displacement only as a sign/order-of-magnitude diagnostic until the full
survey table, coordinates, reference frame, and covariance are supplied.}
\mqPair{mqBlocked}{Blocked}
{What is still blocked for Geoscope, extensometer, crackmeter, laser, and auxiliary
HM streams?}
{Residual-ready Geoscope time series, registered/statistical laser-scan products,
full levelling survey tables, exact units, coordinates/supports, reference-zero
conventions, sign conventions, quality flags, and uncertainty models are missing.}
\mqSeparator{Additional granular questions for OGS grounding}
\mqPair{mqGate}{Gate}
{What quantity does each extensometer row measure?}
{Expected quantity is displacement or strain along extensometer geometry, but no
numeric Geoscope extensometer export is present.}
\mqPair{mqGate}{Gate}
{What quantity does each convergence row measure?}
{Expected quantity is tunnel/niche convergence or point-pair displacement.  The local
catalogue has support evidence but no residual-ready convergence time series.}
\mqPair{mqGate}{Gate}
{What quantity does each miniprisma row measure?}
{Expected quantity is geodetic point displacement.  The layout has miniprisma support
zones, but no numeric time series or survey covariance.}
\mqPair{mqGate}{Gate}
{What quantity does each laser-scan product measure?}
{Expected quantity is registered surface displacement or scan-difference/statistical
surface change.  Locally, only surface geometry and minutes evidence are present.}
\mqPair{mqGate}{Gate}
{What quantity does each levelling row measure?}
{The slide summary gives vertical height difference, but the full survey export is
still missing.}
\mqPair{mqGate}{Gate}
{What baseline epoch is used by each HM stream?}
{Known only for levelling: first measurement on 2022-08-29/30 to the 2026-03
campaign.  Other HM baselines are not in the local exports.}
\mqPair{mqGate}{Gate}
{What sign convention is used by each HM stream?}
{Not available except qualitatively for levelling height differences.  Extensometer,
crackmeter, laser, convergence, and miniprisma sign conventions must be supplied.}
\mqPair{mqGate}{Gate}
{What coordinate frame is used by each HM stream?}
{The Tecplot layout gives source-frame geometry, but accepted OGS-frame coordinates
or survey-frame transforms are not available for hard HM residuals.}
\mqPair{mqGate}{Gate}
{What point, line, pair, or surface support is measured?}
{Layout support exists by role: extensometer lines, convergence points, miniprisma
points, laser surfaces, mini-piezometer support, and fracture/crack geometry.  The
numeric exports needed to bind values to those supports are absent.}
\mqPair{mqGate}{Gate}
{How is each HM support projected to the 2D slice?}
{No accepted projection policy exists.  This must be defined stream by stream before
pressure/deformation rows can be weighted in the 2D model.}
\mqPair{mqGate}{Gate}
{What instrument covariance is available?}
{None for hard HM residuals in the current catalogue.  Levelling only reports a
detectable displacement range of 0.1--0.2 mm at 95 percent confidence.}
\mqPair{mqGate}{Gate}
{What survey-network covariance is available?}
{None.  The full levelling survey table and network covariance/reference-frame
metadata are missing.}
\mqPair{mqGate}{Gate}
{How are BCD-A9/B10 failures after September 2025 handled?}
{The May 2026 minutes identify horizontal/vertical extensometer failures in the
closed niche after September 2025.  Until the numeric export supplies flags, those
post-failure records should be excluded or treated as quality-blocked.}
\bottomrule
\end{longtable}

\mqSubtable{Large Faults And Fracture Geometry}
\begin{longtable}{@{}L{0.18\textwidth} L{0.76\textwidth}@{}}
\toprule
\mqPair{mqGate}{Gate}
{How should large faults differ from visible niche cracks in the model?}
{Large 3D faults and fault zones are structural geometry that may become
permeability/porosity priors, masks, anisotropy/tensor-orientation context, or
scenario fields after projection to the 2D slice.  Visible niche cracks and
crackmeters are local surface-deformation measurements that may constrain relative
displacement, crack aperture, or EDZ priors.}
\mqSeparator{Additional granular questions for OGS grounding}
\mqPair{mqGate}{Gate}
{Which \texttt{.dat} file defines the large fault geometry?}
{\path{VisualisationCDA.dat} is the local Tecplot-style layout source containing
\texttt{Kluft} fracture/crack geometry zones.}
\mqPair{mqGate}{Gate}
{Which Tecplot file defines the \texttt{Kluft} geometry?}
{\path{cda_knowledge_base/measurements_info/coordinates_geometry_layout/source_files/VisualisationCDA.dat}
and the copied equivalents in Taupe/other-HM measurement-info folders.}
\mqPair{mqGate}{Gate}
{What coordinate system is used by the fault source file?}
{The file is in the source layout/Tecplot coordinate frame; the accepted transform
into the current 2D OGS slice is not fully closed for individual faults.}
\mqPair{mqGate}{Gate}
{Does the source contain only geometry?}
{For the local \texttt{Kluft} zones, yes for current model use: the audit extracts
geometry/labels and bounds, not residual-ready aperture or time-series values.}
\mqPair{mqGate}{Gate}
{Does the source contain measured aperture?}
{Not found in the extracted \texttt{Kluft} geometry inventory.  Aperture remains a
separate data request or structural interpretation item.}
\mqPair{mqGate}{Gate}
{Does the source contain measured displacement?}
{Not as a residual-ready fault-displacement table.  Crackmeter/extensometer plots are
separate HM streams and are not the large-fault geometry source.}
\mqPair{mqGate}{Gate}
{Does the source contain hydraulic-response information?}
{Not directly.  Hydraulic effect must be inferred from permeability, ERT, pressure, or
other observations and should not be assigned solely from geometry.}
\mqPair{mqGate}{Gate}
{Which 3D features intersect the current OGS slice?}
{Not determined in an accepted way.  The local inventory has 22 fracture/crack
geometry zones, but a feature-by-feature 3D-to-2D intersection table is still needed.}
\mqPair{mqGate}{Gate}
{What effective 2D width is assigned to each fault feature?}
{None yet.  Width/support assignment is a modelling-policy gate before faults can
modify material fields.}
\mqPair{mqGate}{Gate}
{Should the fault modify permeability?}
{Potentially, as a structural prior, mask, or scenario field for \(\KK\), but only
after projection and independent evidence justify the anomaly.}
\mqPair{mqGate}{Gate}
{Should the fault modify porosity?}
{Potentially as an EDZ/fracture prior, but no direct porosity modification is active
from the fault file.}
\mqPair{mqGate}{Gate}
{Should the fault modify EDZ classification?}
{Yes as a plausible scenario/prior use; public CD-A characterization links fault
zones, facies, EDZ, and water-content/permeability interpretation.}
\mqPair{mqGate}{Gate}
{Should the fault modify tensor orientation?}
{Possibly.  Bedding/fault orientation may inform tensor orientation, but the current
workflow keeps the exact anisotropy-angle release as a structural gate.}
\mqPair{mqGate}{Gate}
{What uncertainty represents 3D-to-2D projection error?}
{No numeric uncertainty is assigned.  Projection error should be represented by a
scenario/prior width or covariance inflation once the feature intersection is chosen.}
\mqPair{mqGate}{Gate}
{Is the fault static over the experiment?}
{Treat as static structural context by default.  No dated aperture/displacement
evolution is present in the current fault-geometry source.}
\mqPair{mqGate}{Gate}
{Can the fault open or close mechanically over time?}
{Physically possible, but not represented by the current OGS material input unless a
time-dependent fracture/EDZ scenario is explicitly added.}
\mqPair{mqGate}{Gate}
{Can the fault self-seal or heal over time?}
{Physically possible for clay/EDZ behaviour, but not contained in
\path{VisualisationCDA.dat}; it belongs to future time-dependent material or missing
physics scenarios.}
\mqPair{mqGate}{Gate}
{What independent evidence justifies a local permeability anomaly?}
{Use permeability pulse tests, ERT high-resistivity/fracture indications, pressure
response, NMR/Taupe trends, or visible crack/HM evidence.  Geometry alone is not
enough for a hard anomaly.}
\bottomrule
\end{longtable}

\mqSubtable{Other Candidate Measurements}
\begin{longtable}{@{}L{0.18\textwidth} L{0.76\textwidth}@{}}
\toprule
\mqSeparator{Additional granular questions for OGS grounding}
\mqPair{mqGate}{Gate}
{Is atmospheric pressure available as a local time series?}
{Not as a residual-ready CD-A auxiliary export in the current catalogue.  Public
seasonal-HM sources show atmospheric/climate monitoring in related Mont Terri niches,
but the project-specific series is not present here.}
\mqPair{mqGate}{Gate}
{Is tunnel temperature available as a local time series?}
{Mentioned in Geoscope/RH context, but no machine-readable local temperature series
is present for boundary filtering or Kelvin conversion.}
\mqPair{mqGate}{Gate}
{Is evapometer conductivity available as a local time series?}
{No local time series was found.  The Tecplot layout includes an evapometer support
role, but not measured values or units.}
\mqPair{mqGate}{Gate}
{Are geophone or MSM velocity data available?}
{Not found as local residual-ready data for the current OGS workflow.  Public CD-A
characterization mentions miniseismic context, but no geophone/MSM export is
catalogued here.}
\mqPair{mqGate}{Gate}
{Are CCM water-content data available?}
{Public/local CD-A sources mention CCM water-content measurements, but the current
processed OGS target set does not include a residual-ready CCM table.}
\mqPair{mqGate}{Gate}
{Are borehole images available for the current support?}
{Not found as a processed measurement stream in the current catalogue.  Structural
PDF/figure evidence exists, but no row-level borehole-image interpretation table is
available.}
\mqPair{mqGate}{Gate}
{Are geological or lithological logs available for the current support?}
{Partly as structural/geology context from Ziefle/GETE material, bedding notes, and
local source figures.  They are prior/context evidence, not a numeric residual.}
\mqPair{mqGate}{Gate}
{Should this stream be an OGS input?}
{Only if it supplies boundary forcing, initial conditions, material fields, or
structural masks with units/support/uncertainty.  RH/temperature/door events are the
most likely input candidates.}
\mqPair{mqGate}{Gate}
{Should this stream be an OGS output residual?}
{Only if it measures an OGS-predicted quantity with support and uncertainty:
pressure, displacement/strain, water content, resistivity through an observation
operator, or surface displacement.}
\mqPair{mqGate}{Gate}
{Should this stream be a material prior?}
{Yes for geology, faults, borehole images, facies, cracks, and lab/field
permeability or porosity evidence when they constrain \(\KK\), \(\por\), EDZ, or
anisotropy.}
\mqPair{mqGate}{Gate}
{Should this stream be a support mask or quality filter?}
{Yes for door openings, failed sensors, bad ERT timesteps/electrodes, outside-mesh
supports, crack/fault masks, and unreliable RH intervals.}
\mqPair{mqGate}{Gate}
{What single missing item blocks use of this stream?}
{Usually a machine-readable export with timestamps/epochs, units, coordinates or
support, reference convention, quality flags, and uncertainty.  For Taupe/ERT/RH the
blocker is instead calibration/support/provenance.}
\mqPair{mqGate}{Gate}
{Who can provide the missing source or decision?}
{The current request packages identify BGR/Gesa or the responsible measurement
provider as the source for Geoscope/HM exports, RH boundary provenance, Taupe
calibration, ERT transform/covariance, and older permeability endpoints.}
\bottomrule
\end{longtable}
\endgroup
